% AA-40 | Audyt I: kontrola dowodów i uzupełnienia o włókna skończone oraz rozgałęzienie.
\chapter{Całkowitość i twierdzenia o ideałach pierwszych}
\label{aa:ch:integral-extensions}

\section{Elementy całkowite i skończoność modułowa}
W całym rozdziale zakładamy, że $A\subseteq B$ są
przemiennymi pierścieniami z jedynką i wspólnym elementem $1$.
Pierścienie te nie muszą być dziedzinami całkowitości, chyba że
zostało to osobno zaznaczone. Rozróżniamy dwa pojęcia:
\emph{element algebraiczny} nad ciałem oraz \emph{element całkowity}
nad pierścieniem. Drugie wymaga równania monicznego, a nie
jakiegokolwiek równania wielomianowego. Różnica jest istotna
nawet wewnątrz tego samego ciała ułamków: każdy element
$\mathbb Q$ jest algebraiczny nad $\mathbb Q$, ale
np. $1/2$ nie jest całkowity nad $\mathbb Z$.
Dalej pokażemy, dlaczego moniczność pozwala sprowadzić
wszystkie wysokie potęgi elementu do skończenie wielu
niższych potęg.

\begin{definicja}[Całkowitość]\label{aa:def:integral-element}
Element $b\in B$ jest \emph{całkowity nad $A$}, jeśli istnieją
$n\ge1$ i $a_0,\ldots,a_{n-1}\in A$ takie, że
\[b^n+a_{n-1}b^{n-1}+\cdots+a_1b+a_0=0.\]
Rozszerzenie $A\subseteq B$ jest \emph{całkowite}, jeżeli każdy
element $B$ jest całkowity nad $A$. Mówimy, że $B$ jest
\emph{skończony modułowo nad $A$}, gdy istnieją
$b_1,\ldots,b_m\in B$ takie, że
$B=Ab_1+\cdots+Ab_m$.
\end{definicja}

\begin{przyklad}[Pierwiastki kwadratowe i ułamki]\label{aa:prz:integral-sqrt}
Element $\sqrt2$ jest całkowity nad $\mathbb Z$, gdyż spełnia
moniczne równanie $T^2-2=0$. Stąd także
$3+2\sqrt2$ jest całkowity, co wynika z ogólnego twierdzenia
poniżej; można sprawdzić to bezpośrednio z równania
$(T-3)^2-8=0$, czyli $T^2-6T+1=0$.
Natomiast $1/2\in\mathbb Q$ jest algebraiczny nad
$\mathbb Q$, ale nie jest całkowity nad $\mathbb Z$.
Gdyby spełniał moniczne równanie nad $\mathbb Z$ stopnia $n$,
po pomnożeniu przez $2^n$ otrzymalibyśmy
$1+2a_{n-1}+\cdots+2^na_0=0$, co jest niemożliwe,
ponieważ lewa strona jest liczbą nieparzystą.
\end{przyklad}

\begin{lemat}[Skończony moduł a równanie moniczne]
\label{aa:lem:integrality-determinant}
Niech $b\in B$. Następujące warunki są równoważne:
\begin{enumerate}[label=\textup{(\roman*)}]
\item element $b$ jest całkowity nad $A$;
\item podpierścień $A[b]$ jest skończenie generowanym $A$-modułem;
\item istnieje skończenie generowany $A$-podmoduł $M\subseteq B$
 zawierający $1$ i spełniający $bM\subseteq M$.
\end{enumerate}
\end{lemat}
\begin{proof}
Załóżmy (i). Jeżeli $b^n$ jest $A$-kombinacją niższych
potęg, to po pomnożeniu przez $b$ to samo dotyczy $b^{n+1}$;
indukcja względem wykładnika pokazuje, że każda potęga $b$
należy do $A$-modułu rozpiętego przez
$1,b,\ldots,b^{n-1}$. Stąd (ii). Warunek (ii) implikuje (iii),
ponieważ wystarczy wybrać $M=A[b]$.

Niech zachodzi (iii), $M=Am_1+\cdots+Am_r$, gdzie $1\in M$.
Dla każdego $i$ zapiszmy
$bm_i=\sum_{j=1}^r a_{ij}m_j$ z $a_{ij}\in A$.
Niech $C=(a_{ij})\in M_r(A)$ oraz niech $I_r$ oznacza
macierz jednostkową stopnia $r$. Układ równań możemy zapisać
$(bI_r-C)\mathbf m=0$, gdzie
$\mathbf m=(m_1,\ldots,m_r)^{\mathsf T}$.
To jest równość wektorów w $M^r$, nie zaś równanie
macierzowe nad ciałem: każdy współczynnik macierzy
z pierścienia $B$ działa na odpowiedniej składowej modułu.
Mnożąc przez macierz dołączoną $\operatorname{adj}(bI-C)$
dostajemy $\det(bI-C)m_i=0$ dla każdego $i$.
Ponieważ $1=\sum_i c_im_i$ dla pewnych $c_i\in A$,
wynika stąd $\det(bI-C)\cdot1=0$ w $B$.
Wyznacznik $\det(TI-C)\in A[T]$ jest moniczny stopnia $r$:
w rozwinięciu wyznacznika składnik $T^r$ pochodzi tylko
z iloczynu elementów diagonalnych, natomiast pozostałe
składniki mają stopień mniejszy. Zatem $b$ spełnia
moniczne równanie i zachodzi (i). Argument nie wymaga,
aby $A$ lub $B$ były dziedzinami: tożsamość macierzy
dołączonej i rozwinięcie wyznacznika obowiązują nad każdym
pierścieniem przemiennym. Istotne jest natomiast,
że $M$ zawiera $1$; inaczej równanie mogłoby
anihilować cały $M$ bez anihilowania jedynki pierścienia $B$.
\end{proof}

\begin{twierdzenie}[Domknięcie na działania i przechodniość]
\label{aa:tw:integrality-transitive}
Niech $A\subseteq B\subseteq C$ będą przemiennymi pierścieniami
z jedynką. Wtedy:
\begin{enumerate}[label=\textup{(\roman*)}]
\item jeżeli $b,c\in B$ są całkowite nad $A$, to
$-b$, $b+c$ i $bc$ są całkowite nad $A$;
\item jeżeli $B/A$ i $C/B$ są rozszerzeniami całkowitymi,
 to $C/A$ także jest całkowite.
\end{enumerate}
\end{twierdzenie}
\begin{proof}
Dla (i) z Lematu~\ref{aa:lem:integrality-determinant} wynika,
że $A[b]$ jest skończenie generowanym $A$-modułem.
Ten sam moniczny wielomian dla $c$ ma współczynniki w
$A\subseteq A[b]$, a więc $c$ jest również całkowity
nad $A[b]$. Stąd $A[b,c]=A[b][c]$ jest skończonym
$A[b]$-modułem. Jeżeli
$u_1,\ldots,u_s$ generują $A[b]$ nad $A$, a
$v_1,\ldots,v_t$ generują $A[b,c]$ nad $A[b]$, to
$u_iv_j$ (dla wszystkich $i,j$) generują $A[b,c]$ nad $A$:
współczynniki względem $v_j$ rozwijamy względem $u_i$.
Stosując kryterium (iii) lematu do
$M=A[b,c]$ i do elementów $-b,b+c,bc$ otrzymujemy
ich całkowitość nad $A$.

Dla (ii) weźmy dowolne $z\in C$. Istnieje równanie moniczne
$z^n+b_{n-1}z^{n-1}+\cdots+b_0=0$ z $b_i\in B$.
Każdy $b_i$ jest całkowity nad $A$, więc
$A'=A[b_0,\ldots,b_{n-1}]$ jest skończenie generowanym
$A$-modułem (stosujemy poprzedni argument indukcyjnie).
Z kolei $A'[z]$ jest skończenie generowanym $A'$-modułem.
Przez przechodniość skończoności modułowej $A'[z]$ jest
skończonym $A$-modułem, a lemat pokazuje, że $z$ jest
całkowity nad $A$. Ponieważ $z$ był dowolny, dostajemy (ii).
Zwróćmy uwagę na rolę skończonej liczby współczynników
$b_0,\ldots,b_{n-1}$: nie zakładamy, że cały pierścień
$B$ jest skończonym $A$-modułem. Każdy badany element
$z$ pozwala jednak znaleźć skończoną podalgebrę
$A'\subseteq B$ wystarczającą dla tego dowodu.
\end{proof}


\begin{twierdzenie}[Całkowitość algebry o skończonej liczbie generatorów]
\label{aa:tw:integral-finite-type-finite-module}
Niech $A\subseteq B$ będą przemiennymi pierścieniami z jedynką.
Jeżeli $B=A[b_1,\ldots,b_m]$ dla pewnej skończonej rodziny
elementów $b_i$ całkowitych nad $A$, to $B$ jest skończenie
generowanym $A$-modułem. Odwrotnie, jeśli $B$ jest
skończenie generowanym $A$-modułem, to rozszerzenie $B/A$
jest całkowite.
\end{twierdzenie}
\begin{proof}
Dla $m=0$ mamy $B=A$, czyli moduł jest generowany
przez $1$. Niech $m\ge1$. Dla każdego $i$ wybierzmy
moniczne równanie stopnia $d_i\ge1$ spełnione przez $b_i$,
ze współczynnikami w $A$. Po dołączeniu kolejnych
elementów równanie to pozostaje prawdziwe, ponieważ
$A\subseteq A[b_1,\ldots,b_{i-1}]$. W rezultacie
$A[b_1,\ldots,b_i]$ jest modułem nad
$A[b_1,\ldots,b_{i-1}]$ generowanym przez
$1,b_i,\ldots,b_i^{d_i-1}$. Wykonując indukcyjnie
rozwinięcie współczynników, dostajemy skończoną rodzinę
$\{b_1^{j_1}\cdots b_m^{j_m}:0\le j_i<d_i\}$
generującą $B$ jako $A$-moduł. Rodzina ta nie musi
być liniowo niezależna.

W drugą stronę dla każdego $b\in B$ moduł $M=B$
jest skończenie generowany nad $A$, zawiera jedynkę
i spełnia $bM\subseteq M$. Zatem warunek (iii) z
Lematu~\ref{aa:lem:integrality-determinant} implikuje,
że $b$ jest całkowity nad $A$. Ponieważ $b$ był dowolny,
całe rozszerzenie $B/A$ jest całkowite.
\end{proof}

\begin{przyklad}[Rozszerzenie całkowite nie musi być skończone modułowo]
\label{aa:prz:integral-infinite-extension}
Niech $K=\mathbb F_2$ i niech $L=\overline{\mathbb F}_2$
będzie algebraicznym domknięciem tego ciała. Każdy element
$\alpha\in L$ jest algebraiczny nad $K$, zatem spełnia
moniczne równanie nad $K$: wystarczy podzielić
wielomian anihilujący przez jego niezerowy współczynnik
wiodący. Rozszerzenie $K\subseteq L$ jest więc całkowite.

Nie jest ono jednak skończone modułowo. Dla każdej liczby
$m\ge1$ wewnątrz $L$ istnieje ciało o $2^m$ elementach,
np. zbiór wszystkich pierwiastków wielomianu
$T^{2^m}-T$; pochodna tego wielomianu wynosi $-1$,
więc w ciele $L$ ma on dokładnie $2^m$ różnych pierwiastków.
Dla $q=2^m$ zachodzi $(u+v)^q=u^q+v^q$ oraz
$(uv)^q=u^qv^q$ (charakterystyka $2$). Jeśli $u,v$ są
pierwiastkami, to są nimi także $u+v$, $-u$ i $uv$.
Dla $u\ne0$ równość $u^q=u$ daje
$(u^{-1})^q=u^{-1}$. Zatem zbiór pierwiastków
jest rzeczywiście podciałem $L$ o $2^m$ elementach.
Ponieważ $m$ jest dowolne, $L$ jest nieskończone.
Gdyby $L$ było generowane nad $\mathbb F_2$ przez
$m_0$ elementów jako moduł, miałoby najwyżej $2^{m_0}$
elementów, gdyż każda kombinacja liniowa ma tylko
$2^{m_0}$ możliwych układów współczynników.
Otrzymujemy sprzeczność. Przykład oddziela hipotezę
\emph{całkowitości} od dodatkowego założenia o skończonej
liczbie generatorów algebry w poprzednim twierdzeniu.
\end{przyklad}

\section{Wznoszenie i nieporównywalność ideałów pierwszych}
\begin{lemat}[Ideał maksymalny i jego przeciwobraz]\label{aa:lem:integral-maximal}
Niech $A\subseteq B$ będzie rozszerzeniem całkowitym,
a $\mathfrak q$ ideałem pierwszym $B$.
Wtedy $\mathfrak q$ jest maksymalny w $B$ wtedy i tylko wtedy,
gdy $\mathfrak p=\mathfrak q\cap A$ jest maksymalny w $A$.
\end{lemat}
\begin{proof}
Przede wszystkim odwzorowanie
$A/\mathfrak p\to B/\mathfrak q$ jest
różnowartościowe: jego jądrem jest
$(\mathfrak q\cap A)/\mathfrak p=(0)$.
Rozszerzenie $A/\mathfrak p\subseteq B/\mathfrak q$ jest
całkowite, ponieważ redukcja równania monicznego modulo
$\mathfrak q$ daje moniczne równanie dla dowolnej klasy
$b+\mathfrak q$ ze współczynnikami w $A/\mathfrak p$.
Oba pierścienie ilorazowe są dziedzinami całkowitości.
Jeżeli $\mathfrak q$ jest maksymalny, $B/\mathfrak q$ jest ciałem.
Dla niezerowego $a\in A/\mathfrak p$ jego odwrotność
$a^{-1}\in B/\mathfrak q$ spełnia moniczne równanie
$a^{-n}+c_{n-1}a^{-(n-1)}+\cdots+c_0=0$ nad $A/\mathfrak p$.
Po pomnożeniu przez $a^{n-1}$ dostajemy
$a^{-1}=-(c_{n-1}+c_{n-2}a+\cdots+c_0a^{n-1})\in A/\mathfrak p$.
A więc $A/\mathfrak p$ jest ciałem i $\mathfrak p$ jest maksymalny.

Odwrotnie, jeśli $A/\mathfrak p$ jest ciałem, to
$B/\mathfrak q$ jest dziedziną całkowitości algebraiczną
nad tym ciałem: każdy jej element spełnia niezerowy
wielomian nad $A/\mathfrak p$.
Dla każdego niezerowego $b\in B/\mathfrak q$ istnieje
niezerowy wielomian anihilujący nad ciałem $A/\mathfrak p$.
Spośród wszystkich monicznych wielomianów anihilujących
$b$ wybierzmy $m_b(T)$ o najmniejszym stopniu; możliwość
unormowania współczynnika wiodącego zapewnia fakt,
że pierścień współczynników jest ciałem.
Ma on niezerowy wyraz wolny: gdyby był zerowy,
to $m_b(T)=T h(T)$ i z $0=b h(b)$ w dziedzinie
$B/\mathfrak q$ wynikałoby $h(b)=0$. Ponieważ $b\ne0$,
przeczyłoby to minimalnemu stopniowi $m_b$.
Z równania $m_b(b)=0$ wynika
$b\cdot(-b^{n-1}-c_{n-1}b^{n-2}-\cdots-c_1)=c_0$;
ponieważ $c_0\ne0$ należy do ciała $A/\mathfrak p$,
$b$ ma odwrotność. Zatem $B/\mathfrak q$ jest ciałem. Pokazaliśmy obie
implikacje osobno, wykorzystując założenie pierwszości
$\mathfrak q$ właśnie po to, by iloraz był dziedziną.
\end{proof}

\begin{twierdzenie}[Twierdzenie o istnieniu ideału leżącego nad danym]\label{aa:tw:integral-lying-over}
Jeśli $A\subseteq B$ jest rozszerzeniem całkowitym, to dla
każdego ideału pierwszego $\mathfrak p\subset A$ istnieje
ideał pierwszy $\mathfrak q\subset B$ spełniający
$\mathfrak q\cap A=\mathfrak p$.
\end{twierdzenie}
\begin{proof}
Ustalmy $\mathfrak p$ i połóżmy $S=A\setminus\mathfrak p$.
Lokalizacja $S^{-1}A=A_{\mathfrak p}$ jest pierścieniem
lokalnym o jedynym ideale maksymalnym
$\mathfrak pA_{\mathfrak p}$, co wynika z opisu jednostek:
$a/s$ jest jednostką dokładnie wtedy, gdy $a\notin\mathfrak p$.
Rozszerzenie $A_{\mathfrak p}\subseteq S^{-1}B$ jest
całkowite. Aby to sprawdzić, weźmy dowolny ułamek
$b/s\in S^{-1}B$ i moniczne równanie
$b^n+a_{n-1}b^{n-1}+\cdots+a_0=0$ w $B$.
Po podzieleniu przez $s^n$ dostajemy
\[
(b/s)^n+(a_{n-1}/s)(b/s)^{n-1}
+\cdots+(a_1/s^{n-1})(b/s)+a_0/s^n=0.
\]
Wszystkie współczynniki tego monicznego równania należą
do $A_{\mathfrak p}$, bo $s\in S$.
Jest to równanie w $S^{-1}B$, niezależnie od tego,
czy $B$ ma dzielniki zera.
Pierścień $S^{-1}B$ jest niezerowy: gdyby $1/1=0/1$
w $S^{-1}B$, istniałby $s\in S$ spełniający $s\cdot1_B=0$,
a więc $s=0$ w $B$ i w $A$ z różnowartościowości
włączenia $A\subseteq B$. Przeczyłoby to $0\notin S$.
Wybierzmy jego ideał maksymalny $\mathfrak n$.
Lemat~\ref{aa:lem:integral-maximal} mówi, że
$\mathfrak n\cap A_{\mathfrak p}$ jest maksymalny, więc
równa się $\mathfrak pA_{\mathfrak p}$.
Ideał pierwszy $\mathfrak q$ będący przeciwobrazem
$\mathfrak n$ względem $B\to S^{-1}B$ ma zatem
$\mathfrak q\cap A=\mathfrak p$: jego przecięcie z $A$
jest przeciwobrazem $\mathfrak n\cap A_{\mathfrak p}
=\mathfrak p A_{\mathfrak p}$ względem
$A\to A_{\mathfrak p}$, a ten przeciwobraz to dokładnie
$\mathfrak p$ na mocy
Twierdzenia~\ref{aa:tw:krull-localization-primes}.
\end{proof}

\begin{twierdzenie}[Wznoszenie łańcuchów ideałów pierwszych]\label{aa:tw:integral-going-up}
Niech $A\subseteq B$ będzie rozszerzeniem całkowitym.
Jeśli $\mathfrak p_0\subsetneq\cdots\subsetneq\mathfrak p_r$
jest łańcuchem ideałów pierwszych $A$, a
$\mathfrak q_0$ ideałem pierwszym $B$ z
$\mathfrak q_0\cap A=\mathfrak p_0$, to istnieją ideały pierwsze
$\mathfrak q_0\subsetneq\cdots\subsetneq\mathfrak q_r$
takie, że $\mathfrak q_i\cap A=\mathfrak p_i$ dla każdego $i$.
\end{twierdzenie}
\begin{proof}
Wystarcza podnieść jeden krok; następnie powtórzymy
konstrukcję $r$ razy. Jeśli $\mathfrak q_i$ leży
nad $\mathfrak p_i$, to odwzorowanie
$A/\mathfrak p_i\to B/\mathfrak q_i$ jest
różnowartościowe, ponieważ
$\mathfrak q_i\cap A=\mathfrak p_i$.
Rozszerzenie to jest całkowite: każda klasa elementu $b\in B$
spełnia równanie moniczne otrzymane przez redukcję modulo
$\mathfrak q_i$ równania dla $b$ nad $A$.
Ideał $\mathfrak p_{i+1}/\mathfrak p_i$ jest pierwszy
w $A/\mathfrak p_i$. Twierdzenie~\ref{aa:tw:integral-lying-over}
daje nad nim ideał pierwszy $\overline{\mathfrak q}_{i+1}$
w $B/\mathfrak q_i$. Jego przeciwobraz $\mathfrak q_{i+1}$
w $B$ zawiera $\mathfrak q_i$ i leży nad
$\mathfrak p_{i+1}$. Inkluzja jest ścisła, gdyż kontrakcje
obu ideałów są różne. Powtarzamy to skończenie wiele razy.
\end{proof}

\begin{twierdzenie}[Nieporównywalność ideałów nad tym samym]\label{aa:tw:integral-incomparability}
Niech $A\subseteq B$ będzie rozszerzeniem całkowitym,
a $\mathfrak q_1\subseteq\mathfrak q_2$ ideałami pierwszymi $B$.
Jeśli $\mathfrak q_1\cap A=\mathfrak q_2\cap A$, to
$\mathfrak q_1=\mathfrak q_2$.
\end{twierdzenie}
\begin{proof}
Załóżmy przeciwnie, że inkluzja jest ścisła.
Przechodzimy do całkowitego rozszerzenia dziedzin
$A/\mathfrak p\subseteq B/\mathfrak q_1$: włączenie
wynika z równości kontrakcji $\mathfrak q_1\cap A
=\mathfrak p$, a całkowitość przechodzi na ilorazy, gdzie
$\mathfrak p=\mathfrak q_i\cap A$.
Niezerowy ideał pierwszy $\mathfrak q_2/\mathfrak q_1$
ma przecięcie $0$ z $A/\mathfrak p$.
Wybierzmy w nim niezerowy element $b$. Ponieważ $b$ jest
całkowity, spełnia moniczne równanie minimalnego możliwego
stopnia nad $A/\mathfrak p$:
$b^n+a_{n-1}b^{n-1}+\cdots+a_0=0$.
Ponieważ $b$ należy do ideału, również $a_0$ do niego należy;
a zatem $a_0=0$, gdyż przecięcie z $A/\mathfrak p$ jest zerowe.
Dzieląc równanie przez $b$ w dziedzinie $B/\mathfrak q_1$
(znaczy: skracając niezerowy czynnik $b$), dostajemy moniczne
równanie stopnia $n-1$ dla $b$, wbrew minimalności $n$.
\end{proof}

\begin{twierdzenie}[Skończoność włókien przy rozszerzeniu skończonym modułowo]
\label{aa:tw:integral-finite-fibers}
Niech $A\subseteq B$ będzie rozszerzeniem pierścieni przemiennych
z jedynką, w którym $B$ jest generowany jako $A$-moduł przez
$m\ge1$ elementów. Dla każdego ideału pierwszego
$\mathfrak p\subset A$ istnieje co najmniej jeden, a co najwyżej
$m$ ideałów pierwszych $\mathfrak q\subset B$ spełniających
$\mathfrak q\cap A=\mathfrak p$.
\end{twierdzenie}
\begin{proof}
Z Twierdzenia~\ref{aa:tw:integral-finite-type-finite-module}
rozszerzenie jest całkowite, więc istnienie przynajmniej
jednego takiego ideału wynika z
Twierdzenia~\ref{aa:tw:integral-lying-over}.
Ustalmy $\mathfrak p$ i połóżmy
$S=A\setminus\mathfrak p$ oraz
$\kappa(\mathfrak p)=\operatorname{Frac}(A/\mathfrak p)$.
Rozważmy pierścień
\[
 C=B\otimes_A\kappa(\mathfrak p)
  \cong S^{-1}(B/\mathfrak p B).
\]
Pierwsza postać ilustruje zmianę pierścienia współczynników;
druga wynika z faktu, że lokalizacja i przejście do ilorazu
komutują. Ponieważ $B$ ma $m$ generatorów jako $A$-moduł,
ich obrazy generują $C$ jako przestrzeń liniową nad
$\kappa(\mathfrak p)$. Zatem
$\dim_{\kappa(\mathfrak p)}C\le m$.

Ideały pierwsze $C$ odpowiadają dokładnie ideałom pierwszym
$\mathfrak q$ pierścienia $B$, które zawierają
$\mathfrak p B$ i są rozłączne z $S$. Warunki te są równoważne
równości $\mathfrak q\cap A=\mathfrak p$.
Każdy ideał pierwszy $\mathfrak n$ pierścienia $C$ jest
maksymalny: iloraz $C/\mathfrak n$ jest skończenie wymiarową
dziedziną nad $\kappa(\mathfrak p)$, a więc ciałem.
Istotnie, mnożenie przez dowolny niezerowy element tej
dziedziny jest różnowartościowym endomorfizmem jej
skończenie wymiarowej przestrzeni liniowej, zatem jest
surjektywne i element ma odwrotność.

Gdy $\mathfrak n_1,\ldots,\mathfrak n_r$ są parami różnymi
ideałami maksymalnymi $C$, są parami współmaksymalne.
Chińskie twierdzenie o resztach daje surjekcję
$C\twoheadrightarrow\prod_{j=1}^r C/\mathfrak n_j$.
Każdy czynnik po prawej ma wymiar nad
$\kappa(\mathfrak p)$ co najmniej $1$, więc
$r\le\dim_{\kappa(\mathfrak p)}C\le m$.
Wobec dowolności tej rodziny istnieje co najwyżej
$m$ ideałów maksymalnych $C$, a zatem tyle samo
ideałów pierwszych $B$ leżących nad $\mathfrak p$.
\end{proof}

\begin{uwaga}\label{aa:uw:integral-fibers-finite-module}
Ograniczenie liczby ideałów w powyższym twierdzeniu
korzysta ze \emph{skończoności modułowej}, nie tylko
z całkowitości. Natomiast twierdzenia o istnieniu
ideału leżącego nad danym, wznoszeniu i nieporównywalności
wymagają jedynie całkowitości rozszerzenia.
\end{uwaga}

\begin{wniosek}[Wymiar nie zmienia się w rozszerzeniu całkowitym]\label{aa:cor:integral-dimension}
Jeśli $A\subseteq B$ jest rozszerzeniem całkowitym niezerowych
pierścieni przemiennych z jedynką, to $\dim A=\dim B$.
\end{wniosek}
\begin{proof}
Każdy łańcuch ideałów pierwszych $A$ podnosimy do łańcucha
tej samej długości w $B$ za pomocą
Twierdzeń~\ref{aa:tw:integral-lying-over} i
\ref{aa:tw:integral-going-up}, stąd $\dim A\le\dim B$.
Z kolei weźmy dowolny ścisły łańcuch
$\mathfrak q_0\subsetneq\cdots\subsetneq\mathfrak q_r$
w $B$. Jego kontrakcje $\mathfrak p_i=\mathfrak q_i\cap A$
są ideałami pierwszymi i spełniają
$\mathfrak p_i\subseteq\mathfrak p_{i+1}$.
Żadne dwa sąsiednie ideały nie mogą mieć tej samej kontrakcji,
gdyż przeczyłoby to
Twierdzeniu~\ref{aa:tw:integral-incomparability}.
Stąd $\dim B\le\dim A$.
\end{proof}


\section{Przykłady wznoszenia oraz rola całkowitości}
\begin{przyklad}[Dwa ideały nad $(5)$ w pierścieniu Gaussa]
\label{aa:prz:integral-gaussian-five}
Niech $A=\mathbb Z$ i $B=\mathbb Z[i]$, gdzie $i^2=-1$.
Każdy element $a+bi\in B$ jest całkowity nad $\mathbb Z$:
ma równanie
\[
T^2-2aT+(a^2+b^2)=0,
\]
co sprawdzamy po podstawieniu $T=a+bi$. Zatem $B/A$
jest rozszerzeniem całkowitym.

Rozpatrzmy dwa homomorfizmy
$\varphi_2,\varphi_3:B\to\mathbb F_5$ określone przez
$\varphi_2(a+bi)=\overline a+2\overline b$ oraz
$\varphi_3(a+bi)=\overline a+3\overline b$.
Oba są homomorfizmami pierścieni, ponieważ
$2^2=3^2=-1$ w $\mathbb F_5$.
Są surjektywne: każdą klasę z $\mathbb F_5$ można otrzymać
z odpowiedniej liczby całkowitej. Ich jądra
\[
\mathfrak q_2=(5,i-2),\qquad
\mathfrak q_3=(5,i-3)
\]
są więc ideałami maksymalnymi $B$. Ich kontrakcje do
$\mathbb Z$ wynoszą $(5)$, bo ograniczenie obu
homomorfizmów do $\mathbb Z$ jest redukcją modulo $5$.
Są to dwa różne ideały: element $i-2$ należy do
$\mathfrak q_2$, lecz
$\varphi_3(i-2)=\overline 1\ne0$.
Nie mogą się zawierać jeden w drugim, gdyż oba są
maksymalne. Otrzymujemy konkretny obraz tego, że nad jednym
ideałem pierwszym może leżeć więcej niż jeden ideał pierwszy,
a twierdzenie o nieporównywalności nie orzeka ich jedyności.
\end{przyklad}



\begin{przyklad}[Nierozkładalność nad liczbą pierwszą a liczba ideałów leżących nad nią]
\label{aa:prz:integral-gaussian-three}
Dla $A=\mathbb Z$ i $B=\mathbb Z[i]$ porównajmy
Przykład~\ref{aa:prz:integral-gaussian-five} z ideałem $(3)$.
Rozważmy homomorfizm $\mathbb Z[T]\to B$,
$T\mapsto i$. Jest on surjektywny, a jego jądro jest
ideałem $(T^2+1)$: dzieląc dowolny wielomian przez
moniczny wielomian $T^2+1$, otrzymujemy resztę $a+bT$;
jeżeli $a+bi=0$ w $B\subseteq\mathbb C$, to $a=b=0$.
Zatem
\[
 B/(3)\cong\mathbb F_3[T]/(T^2+1).
\]
W ciele $\mathbb F_3$ wartości $T^2+1$ dla
$T=0,1,2$ są odpowiednio $1,2,2$; wielomian stopnia
$2$ nie ma więc pierwiastka i jest nierozkładalny.
W konsekwencji $B/(3)$ jest ciałem dziewięcioelementowym,
a $(3)\subseteq B$ jest ideałem maksymalnym.
Jeżeli ideał pierwszy $\mathfrak q$ pierścienia $B$ ma
kontrakcję $(3)$ do $\mathbb Z$, to zawiera $(3)$.
Odpowiada więc ideałowi pierwszemu ciała $B/(3)$;
tym ideałem może być wyłącznie $(0)$. Stąd
$\mathfrak q=(3)$ i nad $(3)$ leży dokładnie jeden
ideał pierwszy $B$.
Natomiast nad $(5)$ leżą dwa różne ideały maksymalne,
opisane w poprzednim przykładzie. Ta różnica pokazuje,
że twierdzenie o istnieniu ideału leżącego nad danym
nie orzeka liczby takich ideałów.
\end{przyklad}
\begin{przyklad}[Rozgałęzienie ideału $(2)$ w pierścieniu Gaussa]
\label{aa:prz:integral-gaussian-two}
Dla $B=\mathbb Z[i]$ mamy
\[
 B/(2)\cong\mathbb F_2[T]/(T^2+1)
 =\mathbb F_2[T]/((T+1)^2).
\]
Niech $\varepsilon$ oznacza klasę $T+1$ w ostatnim ilorazie.
Wtedy $\varepsilon^2=0$, ale $\varepsilon\ne0$.
Każdy element tego pierścienia ma jednoznaczną postać
$a+b\varepsilon$ dla $a,b\in\mathbb F_2$;
nieodwracalne są dokładnie elementy ideału
$(\varepsilon)$, gdyż dla $a=1$ zachodzi
$(1+b\varepsilon)^{-1}=1+b\varepsilon$.
Pierścień $B/(2)$ ma zatem dokładnie jeden ideał pierwszy,
równy $(\varepsilon)$, będący jego jedynym ideałem maksymalnym.
Jego przeciwobrazem w $B$ jest
$\mathfrak q=(2,i+1)=(1+i)$; równość wynika z
$2=-i(1+i)^2$ oraz $i+1=1+i$.
Ponadto $(2)=\mathfrak q^2$ jako ideały $B$.

Porównując z Przykładami~\ref{aa:prz:integral-gaussian-three}
i~\ref{aa:prz:integral-gaussian-five}, widzimy trzy
możliwe zachowania liczby pierwszej w tym samym
rozszerzeniu skończonym modułowo:
$(2)$ daje pojedynczy ideał z rozgałęzieniem,
$(3)$ pozostaje ideałem pierwszym, natomiast
$(5)$ ma dwa różne ideały nad sobą.
Nie utożsamiamy przy tym ideału $(2)\subseteq B$
z pierwszym ideałem $\mathfrak q$ leżącym nad $(2)$.
\end{przyklad}

\begin{przyklad}[Bez całkowitości nie ma wznoszenia]
\label{aa:prz:integral-localization-failure}
Włączenie $\mathbb Z\subseteq\mathbb Z[1/2]$ nie jest
całkowite. Rzeczywiście, gdyby $1/2$ spełniał
moniczne równanie o współczynnikach całkowitych,
po przemnożeniu przez odpowiednią potęgę $2$
uzyskalibyśmy sprzeczność między liczbą nieparzystą
i zerem, jak w Przykładzie~\ref{aa:prz:integral-sqrt}.
W pierścieniu $\mathbb Z[1/2]$ liczba $2$ jest jednostką,
gdyż ma odwrotność $1/2$. Żaden ideał właściwy tego
pierścienia nie zawiera zatem $2$. Gdyby istniał
ideał pierwszy $\mathfrak q\subseteq\mathbb Z[1/2]$
z kontrakcją $(2)\subset\mathbb Z$, to $2\in\mathfrak q$,
sprzeczność. Bez założenia o całkowitości nawet twierdzenie
o istnieniu ideału leżącego nad ustalonym może zawieść.
\end{przyklad}

\section{Wymiar algebr skończenie generowanych}
\begin{twierdzenie}[Wymiar pierścienia wielomianów]\label{aa:tw:polynomial-krull-dimension}
Dla każdego ciała $K$ i każdej liczby całkowitej $n\ge0$
zachodzi $\dim K[X_1,\ldots,X_n]=n$.
\end{twierdzenie}
\begin{proof}
Dla $n=0$ otrzymujemy ciało $K$, więc wymiar wynosi $0$.
Dla $n=1$ wynik wynika z
Twierdzenia~\ref{aa:tw:krull-pid}, gdyż $K[X_1]$
jest dziedziną ideałów głównych. Załóżmy, że $n\ge2$
i znamy wynik dla $n-1$ zmiennych.
Łańcuch
$(0)\subsetneq(X_1)\subsetneq(X_1,X_2)\subsetneq\cdots
\subsetneq(X_1,\ldots,X_n)$
składa się z ideałów pierwszych: odpowiednie ilorazy
są pierścieniami wielomianów, czyli dziedzinami.
Wynika stąd $\dim K[X_1,\ldots,X_n]\ge n$.

Dla nierówności odwrotnej weźmy dowolny łańcuch
$\mathfrak p_0\subsetneq\mathfrak p_1\subsetneq\cdots
\subsetneq\mathfrak p_r$. Jeśli $r=0$, nie ma czego dowodzić.
Ideał $\mathfrak p_1$ jest niezerowy, gdyż
$\mathfrak p_0\supseteq(0)$ i inkluzja jest ścisła.
Wybierzmy $0\ne f\in\mathfrak p_1$. Wielomian $f$
nie może być niezerową stałą: stałe niezerowe są jednostkami
w $K[X_1,\ldots,X_n]$, a właściwy ideał pierwszy
$\mathfrak p_1$ nie zawiera jednostki. Dlatego $f$
jest niestały. Stosujemy trójkątną zmianę zmiennych z
Lematu~\ref{aa:lem:noether-monic}. Jest ona automorfizmem:
jeśli $Y_i=X_i-X_n^{m_i}$ dla $i<n$ i $Y_n=X_n$, to
$X_i=Y_i+Y_n^{m_i}$ i $X_n=Y_n$ wyznaczają odwrotną
substytucję. Po tej zmianie niezerowy wielomian $f$
przyjmuje postać moniczną w $Y_n$ (po pomnożeniu przez
niezerową stałą z $K$, co nie zmienia ideału $(f)$). Wobec tego pierścień
$K[X_1,\ldots,X_n]/(f)$ jest skończenie generowanym
modułem nad pierścieniem $K[Y_1,\ldots,Y_{n-1}]$:
klasy $1,Y_n,\ldots,Y_n^{d-1}$ generują go, gdzie $d$
to stopień otrzymanego wielomianu względem $Y_n$.
Włączenie tego pierścienia wielomianów jest różnowartościowe.
Istotnie, jeśli $0\ne h\in K[Y_1,\ldots,Y_{n-1}]$
należałby do ideału $(f)$, to $h=fq$ dla pewnego
$q\in K[Y_1,\ldots,Y_n]$. Jeżeli $q\ne0$, moniczność
$f$ w $Y_n$ daje $\deg_{Y_n}(fq)=d+\deg_{Y_n}q
\ge d>0$, wbrew $\deg_{Y_n}h=0$; jeśli $q=0$,
to $h=0$, także sprzeczność.
Z Lematu~\ref{aa:lem:integrality-determinant} wynika
całkowitość rozszerzenia; Wniosek~\ref{aa:cor:integral-dimension}
i założenie indukcyjne dają
$\dim(K[X_1,\ldots,X_n]/(f))=n-1$.
Łańcuch $\mathfrak p_1\subsetneq\cdots\subsetneq\mathfrak p_r$
po przejściu do ilorazu przez $(f)$ ma długość $r-1$.
Stąd $r-1\le n-1$, czyli $r\le n$.
Ponieważ łańcuch był dowolny, dowód jest zakończony.
\end{proof}

\begin{wniosek}[Wymiar i stopień transcendencji]\label{aa:cor:finite-type-dimension}
Jeżeli $A$ jest niezerową skończenie generowaną przemienną
$K$-algebrą z jedynką, to istnieje $r\ge0$ takie, że
$\dim A=r$ i $A$ jest skończonym modułem nad
podalgebrą $K[y_1,\ldots,y_r]$ izomorficzną z
pierścieniem wielomianów. Jeżeli dodatkowo $A$ jest dziedziną
całkowitości, to
$\dim A=\operatorname{trdeg}_K\operatorname{Frac}(A)$.
\end{wniosek}
\begin{proof}
Lemat normalizacyjny Noether
(Twierdzenie~\ref{aa:tw:noether-normalization}) dostarcza
algebraicznie niezależnych $y_1,\ldots,y_r\in A$,
nad którymi $A$ jest skończenie generowanym modułem.
Włączenie $K[y_1,\ldots,y_r]\subseteq A$ jest więc
całkowite na mocy Lematu~\ref{aa:lem:integrality-determinant}:
każdy element $a\in A$ mnoży skończony moduł $A$ w siebie,
a $1\in A$. Wniosek~\ref{aa:cor:integral-dimension} daje
$\dim A=\dim K[y_1,\ldots,y_r]=r$.
Jeśli $A$ jest dziedziną, Wniosek~\ref{aa:wn:noether-dimension}
identyfikuje $r$ ze stopniem transcendencji jej ciała ułamków.
\end{proof}

\begin{przyklad}[Krzywa o osobliwości ostrzowej]\label{aa:prz:integral-cusp}
Niech $A=K[X,Y]/(Y^2-X^3)$, gdzie $K$ jest dowolnym ciałem.
Oznaczmy klasy zmiennych przez $x,y$. Z równania
$y^2=x^3$ otrzymujemy $A=K[x]+K[x]y$:
każdą wyższą potęgę $y$ redukujemy, zastępując $y^2$ przez $x^3$.
Element $x$ jest transcendentalny nad $K$, gdyż żaden
niezerowy wielomian w $K[X]$ nie należy do ideału
$(Y^2-X^3)$ (jego generator ma dodatni stopień względem $Y$
i jest moniczny). Wobec tego $K[x]\cong K[T]$.
Całkowitość $A/K[x]$ daje $\dim A=1$, chociaż $A$ ma dwa
generatory jako $K$-algebra. Warto porównać ten przykład z
Rozdziałem~\ref{aa:ch:noether-normalization}.
\end{przyklad}



\begin{przyklad}[Normalizacja krzywej ostrzowej — pełne obliczenie]
\label{aa:prz:integral-cusp-normalization}
Niech $K$ będzie dowolnym ciałem, $t$ zmienną,
$A=K[t^2,t^3]\subseteq B=K[t]$. Element $t\in B$
spełnia moniczne równanie $U^2-t^2=0$ o współczynnikach
w $A$. Co więcej, każdy wielomian $h(t)\in K[t]$
rozkłada się na sumę swoich składników o parzystych
i nieparzystych wykładnikach:
$h(t)=a(t^2)+t b(t^2)$ dla pewnych $a,b\in K[U]$.
Wobec $K[t^2]\subseteq A$ oznacza to
$B=A\cdot1+A\cdot t$; zatem $B$ jest skończenie
generowanym $A$-modułem, a rozszerzenie jest całkowite.

Sprawdźmy precyzyjnie związek z równaniem ostrza.
Homomorfizm $K[X,Y]\to K[t]$ określony wzorami
$X\mapsto t^2$, $Y\mapsto t^3$ ma obraz $A$.
Wielomian $Y^2-X^3$ leży w jego jądrze. Dzięki temu,
że $Y^2-X^3$ jest moniczny względem $Y$, każdy
wielomian modulo $(Y^2-X^3)$ można przedstawić jako
$a(X)+Yb(X)$ dla $a,b\in K[X]$. Jeśli jego obraz
$a(t^2)+t^3b(t^2)$ jest zerowy, to oba składniki
muszą znikać oddzielnie: pierwszy zawiera wyłącznie
parzyste potęgi $t$, a drugi wyłącznie nieparzyste
potęgi $t$ (od stopnia $3$). Stąd $a=b=0$.
Jądro homomorfizmu jest zatem \emph{dokładnie}
$(Y^2-X^3)$, a więc
\[
 K[X,Y]/(Y^2-X^3)\cong K[t^2,t^3]=A.
\]
W szczególności $A$ jest dziedziną całkowitości.
Warto zauważyć, że $A$ nie jest izomorficzny z
pierścieniem wielomianów $K[U]$. Istotnie, ideał
$\mathfrak m=(t^2,t^3)$ pierścienia $A$ nie jest główny.
Dla niezerowego $h(t)\in K[t]$ niech
$\nu(h)$ oznacza najmniejszy wykładnik potęgi $t$
występującej w $h$ z niezerowym współczynnikiem.
Wtedy $\nu(hk)=\nu(h)+\nu(k)$, a niezerowe elementy
$A$ mają wartość $\nu$ równą $0$ lub co najmniej $2$.
Gdyby $\mathfrak m=(h)$ dla pewnego $0\ne h\in A$,
z $h\in\mathfrak m$ wynikałoby $\nu(h)\ge2$,
a z $t^2=ha$ dla pewnego $a\in A$ mielibyśmy
$\nu(h)=2$ i $\nu(a)=0$.
Jednocześnie $t^3=hb$ dla pewnego $b\in A$ wymuszałoby
$\nu(b)=1$, co jest niemożliwe.
Ponieważ każdy ideał $K[U]$ jest główny,
$A$ nie może być izomorficzny z $K[U]$.

W pierścieniu $B$ mamy ścisły łańcuch ideałów pierwszych
$(0)\subsetneq(t)$. Kontrakcja $(t)\cap A$ jest
ideałem $(t^2,t^3)$ pierścienia $A$: wielomian w $t^2,t^3$
ma wyraz wolny zerowy dokładnie wtedy, gdy należy do
$(t^2,t^3)$: każda dodatnia potęga $t$ występująca w $A$
ma wykładnik parzysty co najmniej $2$ lub nieparzysty
co najmniej $3$, a zatem jest podzielna w $A$ przez
$t^2$ albo $t^3$. Ten ideał jest maksymalny, gdyż
$A/(t^2,t^3)\cong K$. Widzimy więc konkretny łańcuch
$(0)\subsetneq(t^2,t^3)$ w $A$. Z twierdzenia o
niezmienniczości wymiaru w rozszerzeniu całkowitym wynika
$\dim A=\dim K[t]=1$; pierścień $A$ nie jest wprawdzie
pierścieniem wielomianów w jednej zmiennej, ale ma ten sam
wymiar Krulla co $B$.
\end{przyklad}
\begin{przyklad}[Nie należy mylić algebraiczności z całkowitością]
\label{aa:prz:integral-nonexample}
Rozszerzenie ciał $\mathbb Q/\mathbb Q$ jest algebraiczne,
lecz $1/2\in\mathbb Q$ nie jest całkowity nad $\mathbb Z$.
Dlatego z faktu, że $\operatorname{Frac}(B)$ jest rozszerzeniem
algebraicznym $\operatorname{Frac}(A)$, nie wynika całkowitość
$B/A$. Na przykład $\mathbb Z\subset\mathbb Z[1/2]$
nie jest rozszerzeniem całkowitym, mimo że oba pierścienie mają
to samo ciało ułamków $\mathbb Q$.
\end{przyklad}
